\chapter*{List of Symbols}
\addcontentsline{toc}{chapter}{List of Symbols}
\markboth{List of Symbols}{List of Symbols}

In this work, lowercase letters are generally used for scalar quantities and vectors; vectors are distinguished from scalars by an underline. Capital, Latin letters are used for matrices. Few exceptions to this rule are present to maintain the commonly adopted capitals letters for some scalar quantities (the Young's modulus, for instance). A special font is used to identify the work, the energy, and the potentials. An overline is also used when the quantity is evaluated in a special configuration, generally the reference one.

\section*{Operators}

\begin{tabular}{cp{0.8\textwidth}}

    $\Delta$        & finite difference \\
    $d$             & differential and derivative symbol  \\
    $\partial$      & partial derivative symbol \\
    $\dot{a}$       & differentiation of $a$ with respect to time \\
    $a \, '$        & differentiation of $a$ with respect the curvilinear coordinate \\
    $\| \vec{a} \|$ & norm of the vector $\vec{a}$ \\
    $ a^* $         & predicted value of $a$ in a predictor--corrector algorithm \\
    \textsuperscript{t} & transpose operator \\
    $\otimes$       & dyadic product \\
    $\wedge$        & vector product \\
    $ A^{+} $       & Moore-Penrose pseudoinverse of $ A $
\end{tabular}\\